% Transcription — fonds Alexandre Grothendieck, shelfmark 130, batch 3 (pages 41-49).
% Established from the facsimile archives/batches/130/batch-03.pdf (a mirror of
% https://grothendieck.umontpellier.fr/130.pdf), per the skill
% transcribe-grothendieck. The batch file carries no cover sheet: PDF page N
% is page 40+N of the folder (the Montpellier footer prints « 41 » ... « 49 »),
% and the archivists' pencil numbers bottom-left agree wherever legible
% (« 41 » to « 46 », « 48 », « 49 »; none is visible on page 47). The folder's
% last nine pages; the batch is short because the folder is.
%
% Pass: Opus 5.5 (claude-opus-5-5), 2026-09-23 — first pass, unchecked against the pages by a human.
%
% Eight of the nine pages are transcribed. Absent: page 49, a typed page of
% someone else's text (a draft on integration, « n°357 », « - 82 - ») crossed
% by two long strokes, with nothing of his mathematics on it.
%
% What the batch is. One run, undated, continuing batch 2's local heights
% (M_K the places of K, M_K-bounded sets, admissible functions
% φ : X(K̄) × M_K̄ -> ℝ, the presheaf Φ_X of batch 2's last page). Page 41
% opens mid-sentence on the definition of M_K-bounded parts of X(K̄) × M_K̄
% via an affine cover U_i (margin « à vérifier »).
%   41-44  O_X^* -> Φ_X (f ↦ φ_f) and O_X^* ↪ R_X^*; the sheaf of
%          quasi-functions QF_X = Φ_X ⊔_{O_X^*} R_X^* (amalgamated sum)
%          = Φ_X × R_X^* / θ(O_X^*); the exact sequences
%          Φ_X -> QF_X -> Div_X -> 0 (Cartier divisors) and
%          R_X^* -> QF_X -> Φ_X / Im O_X^* -> 0; Φ_X -> QF_X injective
%          (« quasi-functions with zero divisor are the admissible
%          functions »); kernel of R_X^* -> QF_X = constant sheaf U_X of
%          M_K-units of K when X is geometrically integral (via g constant on
%          irreducible components, algebraic over K); a 3×3 diagram of sheaves.
%   46, 48 Global sections for X geometrically integral: the diagram with
%          rows ending in Pic(X) -> H^1(X, Φ_X) and Div(X) -> H^1(X, Φ_X);
%          claim that for X quasi-projective Γ(X, Φ_X/(O_X^*/U_X)) -> Pic(X)
%          is surjective, reduced to very ample classes, then by functoriality
%          to X = P^r_K, ξ = O_X(1). Page 48 stops after that sentence.
%   45, 47 Cancelled by long diagonal strokes, but legible in part: another
%          formulation of admissibility (loc. M_K-bounded φ : X(K) × M_K -> ℝ
%          seen as φ̃ : X(K) -> Div(M_K) with bounded images of M_K-bounded
%          parts; « Cette condition est-elle suffisante ?? »), then K'/K
%          finite. Mathematically, p. 44's sentence continues on p. 46 and
%          p. 46's on p. 48, while p. 45 runs on into p. 47; the edition keeps
%          the archivists' order and says so in notes. Nothing here says
%          which was written first, and the transcription does not guess.
%
% The hand is his fastest: formulas carry, connecting prose often does not,
% and a great many words stay \ill{}. Batch 2 was looked at (pages 36-40)
% for context only; it had no transcription at the time of writing, so the
% notation is fixed here.
%
% Notation fixed for the batch: sheaves he underlines keep the underline
% inside math (\underline{\mathcal{O}}^*_X, \underline{R}^*_X,
% \underline{QF}_X, \underline{\mathrm{Div}}_X, \underline{U}_X); his
% doubly-underlined Φ is \underline{\Phi}_X; K̄ as \bar K. Square brackets
% inside his text mark his own interlinear additions. A pale curved stroke
% crossing the formulas at the foot of p. 41 and the head of p. 42 is not
% read as a symbol.
\documentclass[11pt,a4paper]{article}
% Le préambule commun à toutes les transcriptions du fonds Grothendieck.
%
% Il tient en une idée : **l'incertitude se compose, elle ne se résout pas.**
% Une transcription qui lisse un mot douteux a détruit la seule chose pour
% laquelle on la faisait — un lecteur doit pouvoir savoir, à chaque endroit,
% ce qui était sur le feuillet et ce qui a été ajouté. Les six macros qui
% suivent sont donc l'appareil critique, et non de la décoration.
%
% Les mêmes commandes sont reconnues par `scripts/render.mjs`, qui produit la
% vue de lecture du site. Une macro ajoutée ici doit l'être aussi là-bas,
% faute de quoi le rendu échouera bruyamment — ce qui est voulu : un rendu
% silencieusement faux est pire qu'un rendu absent.

\ProvidesPackage{grothendieck}[2026/08/08 Transcriptions du fonds Grothendieck]

\usepackage{amsmath,amssymb,amsthm}
\usepackage{mathrsfs}
\usepackage{stmaryrd}
% Les diagrammes commutatifs sont partout dans ce fonds : c'est la langue
% même de la géométrie algébrique de Grothendieck, et une transcription qui
% les rendrait en prose ne transcrirait rien.
\usepackage{tikz-cd}
% Français par défaut — c'est la langue des feuillets — mais l'anglais est
% chargé aussi : la traduction et le résumé sont des documents anglais, et la
% typographie française (espaces avant les deux-points) y serait une faute.
% Ces documents basculent par \selectlanguage{english} en tête de corps.
\usepackage[english,french]{babel}
\usepackage{fontspec}
\usepackage{geometry}
\geometry{a4paper,margin=25mm,marginparwidth=28mm}
\usepackage{marginnote}
\usepackage{xcolor}
% ulem, not soul. `soul`'s \st and \ul are built on a character-by-character
% scan that XeLaTeX's font machinery breaks: the document fails with an
% undefined control sequence rather than degrading. ulem does the same two jobs
% and is Unicode-engine safe. `normalem` stops it hijacking \emph, which the
% transcriptions use for Grothendieck's own emphasis.
\usepackage[normalem]{ulem}
\usepackage{hyperref}
\hypersetup{colorlinks=true,linkcolor=black,urlcolor=black,pdfborder={0 0 0}}

% --- Métadonnées du lot -----------------------------------------------------
% Reprises telles quelles par le script de rendu, qui les affiche en tête de
% la vue de lecture. Elles fixent ce que le fichier prétend couvrir : sans
% elles, une transcription flotte sans point d'attache dans un fonds de seize
% mille feuillets.

\newcommand{\folder}[1]{\gdef\grothFolder{#1}}
\newcommand{\batch}[1]{\gdef\grothBatch{#1}}
\newcommand{\leaves}[2]{\gdef\grothFirst{#1}\gdef\grothLast{#2}}
\newcommand{\foldertitle}[1]{\gdef\grothFoldertitle{#1}}
\gdef\grothFolder{?}\gdef\grothBatch{?}\gdef\grothFirst{?}\gdef\grothLast{?}\gdef\grothFoldertitle{}

% --- Le résumé --------------------------------------------------------------
%
% Il ouvre la lecture modernisée, sous le titre « Résumé », et il est composé
% en petit. Ce n'est pas un ornement : c'est le seul passage du document qui
% ne soit pas des mathématiques, et un lecteur doit pouvoir le reconnaître
% avant de l'avoir lu — puis le sauter s'il connaît déjà le sujet, ou s'y
% arrêter s'il ne le connaît pas. Le filet qui le suit marque la reprise du
% corps.

\newenvironment{resume}
  {\small\setlength{\parskip}{0.45em}}
  {\par\medskip\hrule\medskip}

% --- Mots-clés ---------------------------------------------------------------
%
% En anglais, en clôture du résumé de la lecture modernisée : le vocabulaire
% moderne sous lequel chercher ce que les feuillets construisent. Le manifeste
% du site (`npm run manifest`) les extrait pour étiqueter la cote entière —
% cette ligne est la source unique des tags, il n'y a pas de fichier à côté.
\newcommand{\keywords}[1]{\par\smallskip\noindent
  {\footnotesize\textsc{Keywords} --- \textit{#1}}\par}

% --- L'appareil critique ----------------------------------------------------

% Le numéro de feuillet, en marge. C'est l'ancre qui permet de régler un
% désaccord sur une formule en regardant *un* feuillet plutôt qu'une chemise
% de six cents. Le site s'en sert aussi pour tourner les pages du fac-similé
% quand on fait défiler la transcription.
\newcommand{\leaf}[1]{%
  \marginnote{\footnotesize\color{gray!70}\textsf{#1}}%
  \label{leaf:#1}\ignorespaces}

% Illisible : on ne devine pas. Un mot inventé qui se lit comme les autres est
% le pire résultat possible.
\newcommand{\ill}{\textcolor{red!60!black}{[\,\ldots\,]}}

% Lecture douteuse : le mot est donné, et signalé comme incertain.
\newcommand{\uncertain}[1]{\uline{#1}}

% Ajout de l'éditeur — une lettre manquante, un mot sous-entendu.
\newcommand{\add}[1]{\textcolor{blue!50!black}{[#1]}}

% Biffé par Grothendieck lui-même. Ce qu'il a rayé fait partie du document :
% une rature dit souvent où la pensée a changé de direction.
\newcommand{\struck}[1]{\sout{#1}}

% Note du transcripteur, sur l'état matériel du feuillet.
\newcommand{\note}[1]{\par\smallskip{\small\color{gray!60!black}\textit{[#1]}}\par\smallskip}

% Note marginale de Grothendieck, distincte du corps du texte.
\newcommand{\marginal}[1]{\par\smallskip\noindent\hspace{1em}%
  {\small\color{gray!60!black}#1}\par\smallskip}

% --- Titre ------------------------------------------------------------------

\AtBeginDocument{%
  \begin{center}
    {\large\bfseries Fonds Alexandre Grothendieck --- cote n\textsuperscript{o}~\grothFolder}\\[2pt]
    {\itshape\grothFoldertitle}\\[4pt]
    {\small feuillets \grothFirst--\grothLast\ (lot \grothBatch)}\\[2pt]
    {\footnotesize\color{gray!60!black}Transcription établie d'après le fac-similé numérisé
     par l'Université de Montpellier.\\
     Les lectures incertaines sont soulignées, les illisibles notées [\,\ldots\,],
     les ajouts entre crochets.}
  \end{center}
  \vspace{1em}}

\endinput


\folder{130}
\batch{3}
\pages{41}{49}
\dating{[vers 1971]}
\watermark{Édition de démonstration}
\foldertitle{Hauteurs et symboles locaux : notes manuscrites (s.d.).}

\begin{document}

\section{[Quasi-fonctions et diviseurs]}
\note{titre de l'éditeur, entre crochets : les feuillets n'en portent aucun. Dans ce lot, les crochets droits à l'intérieur de son texte signalent ses propres additions interlinéaires.}

\page{41}
\marginal{à vérifier}
et \uncertain{contenues dans une réunion} \ill{} d'images de parties $M_K$-bornées des $U_i(\bar K) \times M_{\bar K}$.
\note{la page commence au milieu d'une phrase, suite de la p.\ 40 ; ces deux lignes sont soulignées, et « à vérifier », souligné, est écrit en oblique dans la marge gauche.}

Admettons ce point. On aura, par l'\emph{exemple précédent}, un \uncertain{homomorphisme} de faisceaux
\[
  \underline{\mathcal{O}}^*_X \xrightarrow{\ f \mapsto \varphi_f\ } \underline{\Phi}_X
\]
\note{au-dessus de la flèche, entre « $f \mapsto$ » et $\varphi_f$, un mot est biffé, illisible. $\underline{\Phi}_X$ rend un $\Phi$ souligné deux fois, sa notation tout au long du lot.}

d'autre part, on a un hom.\ de faisceaux
\[
  \underline{\mathcal{O}}^*_X \hookrightarrow \underline{R}^*_X
\]
\note{le $\mathcal{O}$ surcharge une autre lettre ; l'étiquette de la flèche, elle aussi surchargée, ne se lit pas (peut-être $f \mapsto f$).}

où $\underline{R}^*_X$ est le faisceau des fonctions rationnelles sur $X$ [N.B.\ nous supposons $X$ réduit ou \uncertain{même} séparable$/K$, pour fixer les idées]. Nous \uncertain{voulons} « une généralisation commune des fonctions admissibles et des fonctions rationnelles », et définissons pour cela \emph{le faisceau des quasi-fonctions} comme le faisceau
\[
  \underline{QF}_X = \underline{\Phi}_X \amalg_{\underline{\mathcal{O}}^*_X} \underline{R}^*_X
  \quad \text{(somme amalgamée)}
\]
\[
  \simeq \underline{\Phi}_X \times \underline{R}^*_X / \theta(\underline{\mathcal{O}}^*_X)
\]
\note{le signe de somme amalgamée est surchargé (un premier signe biffé) ; on le rend par $\amalg$.}

\page{42}
où
\[
  \theta : \underline{\mathcal{O}}^*_X \longrightarrow \underline{\Phi}_X \times \underline{R}^*_X
\]
est défini par
\[
  \theta(f) = (-\varphi_f, f^{\ill{}})
\]
\note{l'exposant de $f$ est un pâté d'encre, illisible.}

On a donc deux suites exactes
\[
  \underline{\Phi}_X \longrightarrow \underline{QF}_X \xrightarrow{\ \mathrm{div}\ } \underline{\mathrm{Div}}_X \longrightarrow 0
  \qquad \bigl(\underline{\mathrm{Div}}_X \simeq \underline{R}^*_X / \underline{\mathcal{O}}^*_X \text{, diviseurs de Cartier}\bigr)
\]
\[
  \underline{R}^*_X \xrightarrow{\ f \mapsto \{f\}\ } \underline{QF}_X \longrightarrow \underline{\Phi}_X / \mathrm{Im}\, \underline{\mathcal{O}}^*_X \longrightarrow 0
\]
\note{« $\simeq \underline{R}^*_X/\underline{\mathcal{O}}^*_X$ » est écrit au-dessus de $\underline{\mathrm{Div}}_X$, « diviseurs de Cartier » à droite de la ligne. L'étiquette $f \mapsto \{f\}$ de la seconde suite est une lecture incertaine. Dans la première suite, le $Q$ de $\underline{QF}_X$ surcharge une autre lettre.}

Les [sections sur $X$ \uncertain{du}] noyau de $\underline{\Phi}_X \to \underline{QF}_X$ est formé des fonctions admissibles $\varphi$ telles qu'il existe $f \in \Gamma(X, \underline{\mathcal{O}}^*_X)$ avec $(1, \varphi) = (f, -\varphi_f)$, ce qui exige $f = 1$ et $-\varphi_f(P, v) = v(f(P)) = 0$, d'où $\varphi = 0$, donc $\underline{\Phi}_X \to \underline{QF}_X$ est injectif :
\emph{les quasi-fonctions de diviseur nul sont les fonctions admissibles sur $X$.}
\note{dans $(1, \varphi)$, le $1$ surcharge un autre signe ; l'ordre des composantes n'est pas celui de la définition de $\theta$ ci-dessus, et la page est transcrite telle quelle.}

Les [sections sur $X$ du] noyau de $\underline{R}^*_X \to \underline{QF}_X$ est formé des $f \in \Gamma(\underline{R}^*_X)$ telles qu'il existe $g \in \Gamma(X, \underline{\mathcal{O}}^*_X)$, t.q.\ $(f, 0) = (g, -\varphi_g)$. On doit avoir

\page{43}
$\varphi_g = 0$ i.e.\ $v(g(P)) = 0$ pour tt $P \in X(\bar K)$, $v \in M_{\bar K}$, donc $g(P)$ doit être une $M_{\bar K}$-unité pour tout $P \in X(\bar K)$.

Si $g$ n'est pas constante sur les composantes [irréd.] \struck{connexes} géom.\ de $X$, alors \ill{} aurait $g(X(\bar K)) = \bar K -$ \uncertain{un fini}, donc \struck{on aurait pour $v \in M_{\bar K}$} comme $v$ [$\in M_{\bar K}$] prend une infinité de valeurs, ce \ill{} \ill{} \ill{} \ill{} serait \uncertain{pas} \uncertain{vérifié}. Ainsi $g$ est algébrique sur $K$, et de plus \ill{} \ill{} qui \uncertain{dit} [la $K$-]\uncertain{algèbre} \uncertain{séparable} des \uncertain{sections} \struck{\ill{}} de $\underline{\mathcal{O}}^*_X$ alg.\ sur $K$, \ill{} doit \struck{être} une $M_K$-\emph{unité}.
\note{la prose de ce paragraphe, entre les formules, est d'une écriture très rapide ; seuls les termes mathématiques se lisent sûrement.}

\emph{Dans} le cas particulier où $X$ est géom.\ intègre, on voit donc que le noyau de $\underline{R}^*_X \to \underline{QF}_X$ est le faisceau \struck{\ill{}} constant $U_X$ défini par le groupe $U$ des $M_K$-unités de $K$. On a donc les suites exactes can.
\note{le second terme de « $\underline{R}^*_X \to \underline{QF}_X$ » est tracé comme un $\Phi$ souligné muni d'un astérisque ; le contexte de la p.\ 42 le fait lire $\underline{QF}_X$, sans certitude.}

\page{44}
\[
  0 \to \underline{\Phi}_X \longrightarrow \underline{QF}_X \longrightarrow \underline{\mathrm{Div}}_X \longrightarrow 0
\]
\[
  0 \to \underline{R}^*_X / \underline{U}_X \longrightarrow \underline{QF}_X \longrightarrow \underline{\Phi}_X / \mathrm{Im}\, \underline{\mathcal{O}}^*_X \longrightarrow 0
\]

\uncertain{Qu'on} a \ill{} \ill{} \ill{} \ill{} diagramme commutatif.

\begin{tikzcd}[column sep=small, row sep=small]
 & & 0 & & \\
0 \arrow[r, dashed] & \underline{\Phi}_X / \mathrm{Im}\, \underline{\mathcal{O}}^*_X \arrow[r, dashed, "\simeq"] & \underline{\Phi}_X / \mathrm{Im}\, \underline{\mathcal{O}}^*_X \arrow[u] \arrow[r, dashed] & 0 & \\
0 \arrow[r] & \underline{\Phi}_X \arrow[r] \arrow[u, dashed] & \underline{QF}_X \arrow[r] \arrow[u] & \underline{\mathrm{Div}}_X \arrow[u, dashed, "\wr"] & \\
0 \arrow[r] & \underline{\mathcal{O}}^*_X / \underline{U}_X \arrow[r] \arrow[u] & \underline{R}^*_X / \underline{U}_X \arrow[r, dashed] \arrow[u] & \underline{\mathrm{Div}}_X \arrow[r, dashed] \arrow[u, dashed, "\wr"] & 0 \\
 & 0 \arrow[u] & 0 \arrow[u] & 0 \arrow[u, dashed, "\wr"] &
\end{tikzcd}

\note{diagramme redessiné sur une grille : les flèches pointillées sont en pointillé sur la page, les autres pleines. Sur la ligne du haut, les deux tirets extérieurs portent eux aussi une marque ondulée semblable à celle de $\simeq$, et leurs pointes ne se voient pas nettement.}

Passant aux sections globales, on trouve (pour $X$ géom.\ intègre [ce qui implique que $\underline{R}^*_X / \underline{U}_X$ est constant, donc flasque]), on obtient le diagramme commutatif ; lignes et colonnes exactes
\note{la phrase se poursuit en tête de la p.\ 46, par le diagramme des sections globales ; la p.\ 45, barrée, s'intercale dans l'ordre des archivistes.}

\page{45}
\note{page entière barrée de longs traits obliques, et lisible en partie seulement. Sa dernière phrase se poursuit en tête de la p.\ 47, barrée de même.}
\[
  \varphi : X(K) \times M_K \longrightarrow \mathbb{R}
\]
est dite \emph{loc.}\ $M_K$-\emph{bornée} [ou \emph{admissible}] (si elle est [$M_K$-]bornée sur les parties $M_K$-bornées, \struck{\uncertain{Notons}} [i.e.\ si pour tt $E$ \ill{}] \ill{} [\uncertain{bornée}] \ill{} \ill{}, $P \in X(K)$ \struck{\ill{}} \struck{$X(\bar K) \times X(\bar K)$}, [\struck{$\mathbb{R}$}] $\times M_K$ \ill{} $M_K$-bornée), donc si $\varphi$ est loc.\ $M_K$-bornée, pour $P \in X(K)$, \struck{\ill{}} $\varphi(P, \cdot) : M_K \to \mathbb{R}$ est un \struck{$M_{\bar K}$ -- $M_{\bar K}$-diviseur} $M_K$-diviseur. Donc une $\varphi$ admissible s'interprète aussi comme une application
\[
  \tilde\varphi : X(K) \longrightarrow \mathrm{Div}(M_K)
\]
qui a [tout au moins] la propriété suivante : si $F$ est une partie $M_K$-bornée de $X(K)$ [i.e.\ telle que $F \times M_K$ soit $M_K$-bornée], \emph{alors} l'ensemble des $\tilde\varphi(P)$ ($P \in F$) est \emph{borné} dans l'ensemble des diviseurs de $M_K$ [considéré comme ens.\ ordonné].
\note{le haut de la page est très surchargé : additions interlinéaires, mots biffés, et un cadre tracé à droite autour de « tt $E$ … bornée ». Dans $X(K)$ et $M_K$ de la première ligne, le $K$ est repassé sur une autre lettre, peut-être $\bar K$.}

Cette condition est-elle suffisante ??
\marginal{?}

Une fonction $\varphi : X(K) \times M_K \to \mathbb{R}$ \ill{} est $M_K$-bornée est aussi dite \ill{} [\ill{}] \ill{} \ill{} telle est \uncertain{loc.}\ admissible),
\note{dernière ligne écrite serrée au bord du feuillet, avec une addition interlinéaire illisible.}

\page{46}
\begin{tikzcd}[column sep=small, row sep=small, nodes={font=\scriptsize}]
 & 0 & 0 & 0 & 0 \\
0 \arrow[r] & \Phi(X) / (\mathbf{G}_m(X)/U) \arrow[r] \arrow[u] & \Gamma(X, \underline{\Phi}_X / (\underline{\mathcal{O}}^*_X / \underline{U}_X)) \arrow[r] \arrow[u] & \mathrm{Pic}(X) \arrow[r] \arrow[u] & H^1(X, \underline{\Phi}_X) \arrow[u] \\
0 \arrow[r] & \Phi(X) \arrow[r] \arrow[u] & QF(X) \arrow[r] \arrow[u] & \mathrm{Div}(X) \arrow[r] \arrow[u] & H^1(X, \underline{\Phi}_X) \arrow[u, no head, "\Vert" description] \\
0 \arrow[r] & \mathbf{G}_m(X)/U \arrow[r] \arrow[u] & R(X)^*/U \arrow[r] \arrow[u] & \mathrm{Div}_{\ell}(X) \arrow[r] \arrow[u] & 0 \arrow[u] \\
 & 0 \arrow[u] & 0 \arrow[u] & 0 \arrow[u] &
\end{tikzcd}

\note{sur la ligne du haut, une première flèche de $\mathrm{Pic}(X)$ vers $H^1$ est biffée et récrite. L'indice de $\mathrm{Div}_{\ell}(X)$ est lu $\ell$ sans certitude.}

(utilisant que $H^1(X, \underline{\mathcal{O}}^*_X) \simeq H^1(X, \underline{\mathcal{O}}^*_X / U)$ [$= \mathrm{Pic}(X)$] grâce au fait que $U$ est flasque).

Je dis que si $X$ est \emph{quasi-projective}, alors dans le diagramme commutatif précédent on peut remplacer $H^1(X, \underline{\Phi}_X)$ par $0$ (deux fois), en d'autres termes que
\[
  \Gamma(X, \underline{\Phi}_X / (\underline{\mathcal{O}}^*_X / \underline{U}_X)) \longrightarrow \mathrm{Pic}(X)
\]
est surjectif. Comme tt $\xi \in \mathrm{Pic}(X)$ est la différence de deux éléments très amples, on est ramené à prouver que tt élément très ample est dans l'image, et
\note{« $= \mathrm{Pic}(X)$ » est écrit sous $H^1(X, \underline{\mathcal{O}}^*_X)$. La phrase se poursuit en tête de la p.\ 48 ; la p.\ 47, barrée, s'intercale.}

\page{47}
\note{page barrée de trois longs traits obliques ; elle poursuit la dernière phrase de la p.\ 45.}
cela signifie aussi que $\tilde\varphi(X(K))$ est une partie bornée de $\mathrm{Div}\, M_K$.

Si $K'$ est une ext.\ finie de $K$, alors on a $M_{K'}$, et $M_{K'} \to M_K$. Toute fonction
\[
  \varphi : X(K') \times M_{K'} \longrightarrow \mathbb{R}
\]
définit une fonction
\[
  \varphi : X(K) \times M_K \longrightarrow \mathbb{R}
\]
\note{dans la première formule, l'argument de $X$ est surchargé : $K'$ ou $\bar K$, lecture incertaine. La page s'arrête là.}

\page{48}
par fonctorialité ceci nous ramène au cas où $X = \mathbb{P}^r_K$, $\xi = \underline{\mathcal{O}}_X(1)$.
\note{suite de la dernière phrase de la p.\ 46 ; le feuillet ne porte rien d'autre de sa main.}

\end{document}
